\manualchapter{Envelopes, DAC mode and patches}{sec:operation}
\subsection{Eight envelope shapes}
The mode button cycles through the eight shapes in
Table~\ref{tab:envelope-generator}.
\begin{table}[ht]\centering
\begin{tabularx}{\textwidth}{@{}lX@{}}
\toprule Index & Shape \\\midrule
0 & Rise, then return to zero \\
1 & Fall, then stay at zero \\
2 & Rise, then hold high \\
3 & Fall, then return high \\
4 & Repeating rise \\
5 & Repeating fall \\
6 & Repeating triangle, starting upward \\
7 & Repeating triangle, starting downward \\\bottomrule
\end{tabularx}
\caption{Shared envelope modes.}\label{tab:envelope-generator}
\end{table}

The first is the
reset/default mode. Mode is saved in the patch; randomize chooses one of
the eight shapes. Sync restarts the sequence. Envelope frequency controls
the traversal rate, not a separate attack and release pair.

For a repeating bass, enable Envelope on A, choose repeating fall, and lower
Envelope Frequency. For noise percussion, disable Tone, enable Noise and
Envelope, and send a trigger to envelope sync with a falling one-shot shape.
The internal envelope controls the chip amplitude when selected; it is not
an independent voltage output.
